\chapter{System Architecture and Evolution}

\section{The Early Version}

The original system was much closer to a data-driven player viewer than the current research platform. The early scoring logic (\textit{add\_new\_players.js}) used simple formulas. The early \textit{build\_graph.py} was Python-centred: it built a relationship graph from matches and produced HTML and PyVis exports. It was along this path that cluster profiles eventually became the basis for the Tribal NeuroSim simulation input.

\section{The Current System}

The system is built on several well-separated layers. Its components are: Python (data collection, graph construction, legacy exports), SQLite (central persistence), Node.js/Express (API shell and orchestration), Rust (pathfinding, graph analytics, global artifact export), React/Vite (frontend and visualisation), C\# ASP.NET Core (standalone database browser), and the Tribal NeuroSim simulation engine.

The component diagram can be found in chapter 5 of the extended version \cite{premadegraph-full}; in the lite version the general description is provided in the text.

\begin{table}[H]
\centering
\caption{Rejected approaches and rationale}
\begin{tabularx}{\textwidth}{|>{\raggedright\arraybackslash}p{3.5cm}|>{\raggedright\arraybackslash}p{3.6cm}|X|}
\hline
\textbf{Alternative} & \textbf{Brief advantage} & \textbf{Why it was not chosen} \\
\hline
Everything in Python & Faster initial iteration & Weaker performance and memory control for runtime graph algorithms and larger exports \\
Everything in Rust & Strong typing and a unified engine & Collection and rapid experimentation would have become significantly harder and slower \\
Frontend-centred computation & Simpler demo architecture & Would have pushed compute-intensive graph work to the browser, weakening reproducibility \\
Server-side SQL database & Stronger multi-user model & Given the project's local, demonstrable, research-oriented nature, SQLite remained the better trade-off \\
\hline
\end{tabularx}
\end{table}

\section{Multi-Dataset Support}

The implementation manages multiple independent datasets. Based on the \textit{backend/data/datasets.json} registry, each dataset has its own raw database, refined database, match JSON directory, and cache directory.

Dataset selection does not stop at the Node layer. The backend derives environment variables from the selected dataset for the Rust runtime (\textit{DATASET\_ID}, \textit{GRAPH\_DB\_PATH}, \textit{PATHFINDER\_MATCH\_DIR}, \textit{PATHFINDER\_BIRDSEYE\_CACHE\_DIR}).

\begin{figure}[H]
\centering
\includegraphics[width=0.98\textwidth,height=0.74\textheight,keepaspectratio]{dataset lifecycle riotapi to analytical output.pdf}
\caption{Lifecycle diagram of a dataset from the Riot API to analytical output}
\end{figure}

\section{Current Data State}

The active data state was verified directly from the current repository state on 28 April 2026.

\begin{table}[H]
\centering
\caption{Current data state of the two datasets (28 April 2026)}
\begin{tabularx}{\textwidth}{|>{\raggedright\arraybackslash}p{4.5cm}|r|r|}
\hline
\textbf{Metric} & \textbf{flexset} & \textbf{soloq} \\
\hline
Match JSON files & 4\,981 & 2\,001 \\
Refined players & 15\,421 & 6\,768 \\
Clusters & 599 & 140 \\
Cluster members & 5\,741 & 1\,538 \\
\hline
\end{tabularx}
\end{table}

The soloq dataset clusters typically represent \emph{duo pairs}: since Solo/Duo Queue allows premades of at most two players, the vast majority of clusters that emerge from recurring connections contain only two or three players.

\section{Persistence and the Relational Model}

Clusters previously existed only in JSON export form; database integration made them first-class database objects. The schema is built around three central tables \cite{sqlite}: \textit{players}, \textit{clusters}, and \textit{cluster\_members}.

\begin{figure}[H]
\centering
\includegraphics[width=0.98\textwidth,height=0.76\textheight,keepaspectratio]{sqlite-project-er-diagram.png}
\caption{Entity-relationship diagram of the main SQLite database}
\end{figure}

\begin{figure}[H]
\centering
\includegraphics[width=0.68\textwidth]{relational model diagram.pdf}
\caption{Relational model for cluster and player persistence (\textit{players}, \textit{clusters}, \textit{cluster\_members})}
\end{figure}

Two cluster families coexist in the system. \textit{python\_population} type clusters are the output of the Python graph-building pipeline: NetworkX-based modularity maximisation with a community-detection perspective. \textit{rust\_pathfinding} type clusters come from the Rust engine's runtime clustering: strongly connected segments based on ally connections, capped at 20 members. The aggregated performance metrics of the cluster profiles also form the genome sequence input for the Tribal NeuroSim simulation.
